\documentclass[11pt,a4paper]{article}

% REQUISITOS TRANSVERSALES
\usepackage{booktabs}
\usepackage{listings}
\usepackage{graphicx}
\usepackage{amsmath}
\usepackage[style=ieee,backend=biber,sorting=none]{biblatex}
\usepackage{hyperref}

% IDIOMA Y TIPOGRAFÍA
\usepackage[spanish,es-nodecimaldot]{babel}
\usepackage[T1]{fontenc}
\usepackage[utf8]{inputenc}
\usepackage{lmodern}
\usepackage{microtype}
\usepackage{csquotes}

% DISEÑO GENERAL
\usepackage[
    a4paper,
    top=1.65cm,
    bottom=1.65cm,
    left=1.75cm,
    right=1.75cm,
    headheight=15pt,
    headsep=7pt,
    footskip=16pt
]{geometry}

\usepackage{xcolor}
\usepackage{titlesec}
\usepackage{fancyhdr}
\usepackage{tabularx}
\usepackage{array}
\usepackage{ragged2e}
\usepackage{setspace}
\usepackage{lastpage}
\usepackage[most]{tcolorbox}

\addbibresource{referencias.bib}

% COLORES
\definecolor{azulPrincipal}{HTML}{17365D}
\definecolor{azulSecundario}{HTML}{2F5597}
\definecolor{azulClaro}{HTML}{EEF3F8}
\definecolor{grisClaro}{HTML}{F5F6F7}
\definecolor{grisMedio}{HTML}{D5DAE0}
\definecolor{grisTexto}{HTML}{3D4349}
\definecolor{verdeClaro}{HTML}{EEF6F1}
\definecolor{verdeBorde}{HTML}{3F7656}

% HIPERVÍNCULOS
\hypersetup{
    colorlinks=true,
    linkcolor=azulSecundario,
    citecolor=azulSecundario,
    urlcolor=azulSecundario,
    pdftitle={Análisis técnico de CockroachDB},
    pdfauthor={Jeremy Ruperto Gaibor Rodríguez},
    pdfsubject={Aplicaciones Distribuidas}
}

% TEXTO
\setlength{\parindent}{0pt}
\setlength{\parskip}{3pt}
\setstretch{1.01}
\justifying

% TÍTULOS
\titleformat{\section}
  {\large\bfseries\color{azulPrincipal}}
  {\thesection.}
  {0.5em}
  {}

\titlespacing*{\section}{0pt}{6pt}{2pt}

\newcommand{\sectionrule}{
    \vspace{-1pt}
    {\color{grisMedio}\hrule height 0.45pt}
    \vspace{2pt}
}

% ENCABEZADO Y PIE
\pagestyle{fancy}
\fancyhf{}

\lhead{\small\textcolor{azulPrincipal}{\textbf{FORM-IC-02}}}
\chead{\small Aplicaciones Distribuidas}
\rhead{\small CockroachDB}

\cfoot{
    \small
    Jeremy Ruperto Gaibor Rodríguez
    \hspace{7pt}|\hspace{7pt}
    Página \thepage\ de \pageref{LastPage}
}

\renewcommand{\headrulewidth}{0.45pt}
\renewcommand{\footrulewidth}{0pt}

% TABLAS
\newcolumntype{Y}{>{\RaggedRight\arraybackslash}X}
\renewcommand{\arraystretch}{1.08}

% CAJAS
\newtcolorbox{resumenbox}{
    colback=azulClaro,
    colframe=azulPrincipal,
    boxrule=0.55pt,
    arc=1mm,
    left=3.5mm,
    right=3.5mm,
    top=2mm,
    bottom=2mm,
    before skip=3pt,
    after skip=4pt
}

\newtcolorbox{decisionbox}{
    colback=verdeClaro,
    colframe=verdeBorde,
    boxrule=0.55pt,
    arc=1mm,
    left=3.5mm,
    right=3.5mm,
    top=2mm,
    bottom=2mm,
    before skip=3pt,
    after skip=4pt
}

% LISTINGS
\lstset{
    basicstyle=\ttfamily\small,
    frame=single,
    rulecolor=\color{grisMedio},
    backgroundcolor=\color{grisClaro},
    keywordstyle=\bfseries\color{azulPrincipal},
    breaklines=true,
    showstringspaces=false
}

% BIBLIOGRAFÍA
\defbibheading{bibliography}[\refname]{
    \section*{#1}
    \sectionrule
}

\AtBeginBibliography{
    \small
    \setlength{\bibitemsep}{2pt}
    \setlength{\bibhang}{9pt}
}

\begin{document}

% ENCABEZADO PRINCIPAL
\begin{center}
    {\LARGE\bfseries\color{azulPrincipal}
    Análisis técnico de CockroachDB}

    \vspace{1mm}

    {\normalsize
    Distribución de datos, consistencia CAP y pertinencia para TiendaTech}
\end{center}

\vspace{1mm}
{\color{azulPrincipal}\hrule height 0.8pt}
\vspace{2mm}

\begin{tabularx}{\textwidth}{@{}>{\bfseries\color{azulPrincipal}}l Y
                               >{\bfseries\color{azulPrincipal}}l Y@{}}
\toprule
Autor: & Jeremy Ruperto Gaibor Rodríguez &
Actividad: & FORM-IC-02 \\
Asignatura: & Aplicaciones Distribuidas &
Proyecto: & TiendaTech \\
\bottomrule
\end{tabularx}

\vspace{2mm}

\begin{resumenbox}
\textbf{\textcolor{azulPrincipal}{Síntesis.}}
CockroachDB es una base de datos \textit{Distributed SQL} que conserva
transacciones relacionales mientras distribuye y replica la información
entre varios nodos. Este análisis revisa su organización de datos, su
clasificación según CAP, un caso real y su posible aplicación en TiendaTech.
\end{resumenbox}

% INTRODUCCIÓN
\section{Introducción}
\sectionrule

CockroachDB conserva tablas, consultas SQL y transacciones ACID, pero
distribuye la información entre varios nodos. Para la aplicación se presenta
como una sola base lógica, mientras el motor administra la ubicación y las
réplicas. En despliegues multirregión también permite declarar mediante SQL
la localidad de los datos y los objetivos de supervivencia
\cite{vanbenschoten2022,ajmani2022}.

Su utilidad aparece cuando una aplicación necesita operar en distintas
ubicaciones sin mantener una base independiente por cada sede. En un sistema
pequeño, esa complejidad puede no estar justificada.

% DISTRIBUCIÓN
\section{Distribución de datos}
\sectionrule

Las tablas e índices se representan en un espacio ordenado de claves dividido
en unidades llamadas \textit{rangos}. Estos pueden separarse al crecer,
repartirse entre nodos y rebalancearse cuando cambia la composición del
clúster \cite{distributionlayer,ajmani2022}. Desde la aplicación no es
necesario conocer qué servidor almacena cada fila.

Cada rango mantiene varias réplicas coordinadas mediante Raft. Una escritura
solo se confirma cuando existe acuerdo de un quórum. Con tres réplicas se
requieren dos; si se pierde esa mayoría, el rango afectado deja de aceptar
operaciones que necesitan consenso \cite{replicationlayer,hardrockcase}.

\begin{table}[ht]
\centering
\small
\caption{Modalidades de localidad en CockroachDB}
\begin{tabularx}{\textwidth}{@{}>{\ttfamily}l Y@{}}
\toprule
\textbf{Configuración} & \textbf{Uso principal} \\
\midrule
GLOBAL &
Favorece lecturas desde distintas regiones. \\

REGIONAL BY TABLE &
Mantiene la tabla cerca de una región principal. \\

REGIONAL BY ROW &
Permite asignar cada fila a una región específica. \\
\bottomrule
\end{tabularx}
\end{table}

La modalidad debe elegirse según el patrón de lectura y escritura de cada
tabla, no aplicarse de igual forma a toda la base
\cite{vanbenschoten2022,ajmani2022}.

% CAP
\section{Consistencia según CAP}
\sectionrule

CockroachDB puede analizarse principalmente como un sistema \textbf{CP}.
Durante una partición mantiene la consistencia y la tolerancia a particiones;
el grupo sin mayoría no puede confirmar escrituras. La disponibilidad puede
reducirse temporalmente antes que permitir estados incompatibles
\cite{replicationlayer,transactionlayer}.

La capa transaccional admite operaciones ACID sobre diferentes filas, tablas
y rangos. Bajo aislamiento serializable, las transacciones concurrentes deben
producir un resultado equivalente a una ejecución ordenada
\cite{transactionlayer,vanbenschoten2022}. En TiendaTech, esta propiedad
permitiría registrar un pedido y reducir el inventario dentro de una sola
transacción.

% CASO REAL
\section{Caso real documentado}
\sectionrule

Hard Rock Digital utiliza CockroachDB en una plataforma de apuestas deportivas
que opera en varios estados de Estados Unidos. La solución mantiene una sola
base lógica y combina regiones de AWS, zonas locales y AWS Outposts para
cumplir requisitos de disponibilidad, escalamiento y residencia de datos
\cite{hardrockcase,hardrockblog}.

El caso reporta operación en ocho estados o regiones y cerca de cien nodos
durante periodos de alta demanda. Las fuentes pertenecen a Cockroach Labs, por
lo que son válidas para documentar la implementación, pero no constituyen una
comparación independiente con otras tecnologías.

% PERTINENCIA
\section{Pertinencia para TiendaTech}
\sectionrule

En TiendaTech, productos, marcas y categorías podrían evaluarse como
información global, porque todas las sucursales consultarían el mismo
catálogo. El inventario podría manejarse como regional por fila, asociando
cada existencia con la sucursal responsable
\cite{vanbenschoten2022,ajmani2022}.

Los pedidos y la reducción del stock deberían ejecutarse dentro de una sola
transacción. Así se registraría la venta y se descontaría la existencia, o se
cancelarían ambas operaciones si ocurre un error.

\begin{decisionbox}
\textbf{\textcolor{verdeBorde}{Decisión técnica.}}
Mantendría PostgreSQL y \texttt{postgres\_fdw} en la versión académica actual,
porque permiten demostrar explícitamente la fragmentación, la replicación y
la reconstrucción. Consideraría CockroachDB si TiendaTech operara en varias
sedes reales, con alta concurrencia y necesidad de continuar ante fallos.
\end{decisionbox}

% CONCLUSIÓN
\section{Conclusión}
\sectionrule

CockroachDB distribuye la información mediante rangos replicados y coordinados
por Raft. Su comportamiento es principalmente CP porque exige quórum y
prioriza la consistencia durante una partición. Hard Rock Digital demuestra
que puede utilizarse en un entorno multirregión real.

Para TiendaTech sería una alternativa futura. En la implementación actual,
PostgreSQL sigue siendo suficiente y permite comprender de forma más directa
los mecanismos de distribución exigidos en la asignatura.

% REFERENCIAS
\printbibliography[title={Referencias}]

\end{document}